\documentclass[a4paper,landscape,11pt]{article}

\usepackage{geometry}
\geometry{a4paper,landscape,margin=1cm,top=1cm,bottom=1cm}

\usepackage{amsmath,amssymb}
\usepackage{xcolor}
\usepackage{tikz}
\usetikzlibrary{positioning,arrows.meta,calc,fit,backgrounds,shapes.geometric}

\usepackage[extrafootnotefeatures]{xepersian}

\settextfont[
Scale=1.0,
Extension=.ttf,
Path=/Users/macbook/develope/thesis/thesis_latex/first_edition/styles/fonts/,
BoldFont=XB NiloofarBd,
ItalicFont=XB NiloofarIt,
BoldItalicFont=XB NiloofarBdIt
]{XB Niloofar}

\setdigitfont[
Scale=1.0,
Extension=.ttf,
Path=/Users/macbook/develope/thesis/thesis_latex/first_edition/styles/fonts/,
BoldFont=XB NiloofarBd
]{XB Niloofar}

\pagestyle{empty}
\setlength{\parindent}{0pt}

% ---------- palette ----------
\definecolor{c1}{HTML}{2C5F8A}   % ch1 blue
\definecolor{c2}{HTML}{A8621B}   % ch2 orange
\definecolor{c3}{HTML}{2E6B4F}   % ch3 green
\definecolor{c4}{HTML}{6B3D8A}   % ch4 purple
\definecolor{c5}{HTML}{555555}   % ch5 gray
\definecolor{cred}{HTML}{B02020}
\definecolor{cgy}{HTML}{6A6A6A}
\definecolor{inkk}{HTML}{1A1A1A}

\newcommand{\bul}{\ensuremath{\bullet}~}
\newcommand{\key}[1]{\textbf{#1}}

\tikzset{
  every node/.style={inner sep=5pt, execute at begin node=\setRTL},
  hdr/.style={rounded corners=3pt, draw=#1, line width=1pt, fill=#1!10,
              text width=\colw, align=right, anchor=north, inner sep=6pt},
  arw/.style={-{Stealth[length=3mm,width=2mm]}, line width=0.9pt, draw=black!55},
  arwl/.style={-{Stealth[length=3mm,width=2mm]}, line width=0.9pt, draw=black!35, rounded corners=6pt},
}

\newlength{\colw}

\newcommand{\pagetitle}[2]{%
  \begin{center}
  {\LARGE\bfseries\color{#1} #2}
  \end{center}
  \vspace{-2mm}
}

\begin{document}

% =====================================================================
% PAGE 1 : OVERALL MAP
% =====================================================================
\pagetitle{inkk}{نقشه‌ی روایت پایان‌نامه \raisebox{1pt}{---} در هر فصل چه می‌گویم}

\setlength{\colw}{4.5cm}
\begin{center}
\begin{tikzpicture}[font=\scriptsize]

\node[hdr=c1, minimum height=5.6cm] (a1) at (10.2,0)
 {{\small\bfseries\color{c1} فصل ۱ \raisebox{1pt}{---} مقدمه}\\[3pt]
  {\bfseries پرسش فصل: چرا این مسئله؟}\\[3pt]
  \bul فاز بدون‌کنتراست فقط «مکمل» نیست: در تروما، غربالگری و نارسایی کلیوی \key{تنها تصویر موجود} است.\\[2pt]
  \bul اما دقت بخش‌بندی خودکار روی همین فاز نظام‌مند پایین‌تر است.\\[2pt]
  \bul کم‌دسترس‌ترین بیمار، ضعیف‌ترین خروجی را می‌گیرد \raisebox{1pt}{---} نقطه‌ی شروع کار.};

\node[hdr=c2, minimum height=5.6cm] (a2) at (5.1,0)
 {{\small\bfseries\color{c2} فصل ۲ \raisebox{1pt}{---} کارهای پیشین}\\[3pt]
  {\bfseries پرسش فصل: ادبیات کجا خالی است؟}\\[3pt]
  \bul پنج حوزه مرور می‌شود: یکی مسئله‌ی چتر و چهارتا متناظر چهار زیرمسئله.\\[2pt]
  \bul خروجی فصل یک چیز است و نه بیشتر: \key{چهار جای خالی} که فصل ۳ دقیقاً به همان ترتیب پرشان می‌کند.};

\node[hdr=c3, minimum height=5.6cm] (a3) at (0,0)
 {{\small\bfseries\color{c3} فصل ۳ \raisebox{1pt}{---} روش پیشنهادی}\\[3pt]
  {\bfseries پرسش فصل: چه ساختم؟}\\[3pt]
  \bul سه مؤلفه که یک \key{زنجیره‌ی وابسته}اند و نه سه کار موازی: تشخیص فاز $\leftarrow$ ثبت تصویر $\leftarrow$ سنتز.\\[2pt]
  \bul به‌علاوه‌ی پروتکل ارزیابی، که خودش بخشی از سهم پژوهش است نه فهرست معیار.};

\node[hdr=c4, minimum height=5.6cm] (a4) at (-5.1,0)
 {{\small\bfseries\color{c4} فصل ۴ \raisebox{1pt}{---} آزمایش‌ها}\\[3pt]
  {\bfseries پرسش فصل: چه دیدم؟}\\[3pt]
  \bul یافته‌های مرور مجموعه‌داده‌ها، نتایج ثبت تصویر بین‌فازی، و نتایج سنتز \lr{PLC-Net}.\\[2pt]
  \bul لحن فصل صریح است: هم ادعاهایی که شواهد تأیید کرد و هم \key{دو ادعایی که شواهد رد کرد} گزارش می‌شوند.};

\node[hdr=c5, minimum height=5.6cm] (a5) at (-10.2,0)
 {{\small\bfseries\color{c5} فصل ۵ \raisebox{1pt}{---} نتیجه‌گیری}\\[3pt]
  {\bfseries پرسش فصل: چه ماند؟}\\[3pt]
  \bul جمع‌بندی یافته‌ها، محدودیت‌ها (به‌ویژه دامنه‌ی اعتبار ناشی از سوگیری داده) و کارهای آتی.\\[4pt]
  {\bfseries\color{cred} وضعیت: هنوز نوشته نشده است.}};

\draw[arw] (a1.west) -- (a2.east);
\draw[arw] (a2.west) -- (a3.east);
\draw[arw] (a3.west) -- (a4.east);
\draw[arw] (a4.west) -- (a5.east);

% --- shared spine band ---
\node[draw=black!35, rounded corners=3pt, fill=black!3, text width=25.6cm,
      align=right, anchor=north, inner sep=8pt] (band) at (0,-6.4)
 {{\small\bfseries ستون‌بندی مشترک: یک ترتیب واحد در هر سه فصل}\\[2pt]
  اگر وسط جلسه گم شدید، به این چهار ستون برگردید: هر زیرمسئله را می‌توان مستقل و کامل روایت کرد، با همین ترتیب چهارگانه‌ی مسئله، پیشینه، روش و آزمایش.};

\setlength{\colw}{5.6cm}
\node[hdr=cgy, minimum height=2.4cm] (b1) at (9.45,-8.3)
 {{\bfseries زیرمسئله ۱: مجموعه‌داده‌ها}\\[2pt]
  پیشینه ۲-۳ \ $\leftarrow$\ روش ۳-۲ \ $\leftarrow$\ آزمایش ۴-۱\\[2pt]
  {\color{c3}\bfseries کامل}};
\node[hdr=cgy, minimum height=2.4cm] (b2) at (3.15,-8.3)
 {{\bfseries زیرمسئله ۲: تشخیص فاز}\\[2pt]
  پیشینه ۲-۴ \ $\leftarrow$\ روش ۳-۴ \ $\leftarrow$\ آزمایش: ---\\[2pt]
  {\color{cred}\bfseries نتایج تجربی نوشته نشده}};
\node[hdr=cgy, minimum height=2.4cm] (b3) at (-3.15,-8.3)
 {{\bfseries زیرمسئله ۳: ثبت تصویر}\\[2pt]
  پیشینه ۲-۵ \ $\leftarrow$\ روش ۳-۵ \ $\leftarrow$\ آزمایش ۴-۲\\[2pt]
  {\color{c3}\bfseries کامل \raisebox{1pt}{---} قوی‌ترین بخش تجربی}};
\node[hdr=cgy, minimum height=2.4cm] (b4) at (-9.45,-8.3)
 {{\bfseries زیرمسئله ۴: سنتز \lr{NC\,$\rightarrow$\,CE}}\\[2pt]
  پیشینه ۲-۶ \ $\leftarrow$\ روش ۳-۶ (\lr{PLC-Net}) \ $\leftarrow$\ آزمایش ۴-۳\\[2pt]
  ابلیشن و کنترل‌پذیری \key{آماده}؛ {\color{cred}\bfseries جدول مقایسه و بازوی بخش‌بندی باز}};

\end{tikzpicture}
\end{center}

\vspace{-1mm}
\begin{center}
\fbox{\parbox{25.6cm}{\small\bfseries جمله‌ای که در هر بخش باید تکرار شود:\normalfont\ هدف، بازطراحی مدل بخش‌بندی نیست؛ نشان‌دادن این است که شکاف دقت میان دو فاز تا چه اندازه با مداخله در \textbf{سطح داده و پیش‌پردازش} کاهش می‌یابد.}}
\end{center}

\newpage
% =====================================================================
% PAGE 2 : CHAPTER 1
% =====================================================================
\pagetitle{c1}{فصل ۱ \raisebox{1pt}{---} مقدمه: در هر بخش چه می‌گویم}

\setlength{\colw}{7.7cm}
\begin{center}
\begin{tikzpicture}[font=\footnotesize]

\node[hdr=c1, minimum height=7.4cm] (n1) at (8.55,0)
 {{\small\bfseries\color{c1} ۱-۱ اهمیت مسئله}\\[4pt]
  \bul سی‌تی شکمی چندفازی است: یک اسکن بدون‌کنتراست و یک یا چند اسکن کنتراست‌دار (شریانی/وریدی/تأخیری).\\[3pt]
  \bul تشخیص \lr{HCC} به الگوی افزایش کنتراست و شست‌وشو میان فازها وابسته است.\\[3pt]
  \bul \key{نکته‌ی کلیدی این بخش:} فاز بدون‌کنتراست صرفاً تکمیلی نیست؛ در تروما، غربالگری فرصت‌طلبانه و بیماران با نارسایی کلیوی یا سابقه‌ی حساسیت، تنها تصویر در دسترس است.\\[3pt]
  \bul اما به‌دلیل افت \lr{CNR} میان بافت‌های نرم مجاور، دقت بخش‌بندی روی همین فاز به‌طور نظام‌مند و وابسته به اندام پایین‌تر است.};

\node[hdr=c1, minimum height=7.4cm] (n2) at (0,0)
 {{\small\bfseries\color{c1} ۱-۲ تعریف مسئله}\\[4pt]
  \bul هدف کلی: کاهش شکاف دقت بخش‌بندی چندعضوی میان فاز \lr{NC} و \lr{CE}.\\[3pt]
  \bul دو راهبرد مکمل: \key{هم‌ترازسازی هندسی} فازهای موجود (ثبت تصویر) و \key{بازسازی} فاز غایب (سنتز).\\[3pt]
  \bul اما هر دو به یک زیرساخت مشترک نیاز دارند: شناخت محدودیت داده‌ی موجود، و شناسایی صحیح فاز هر سری.\\[3pt]
  \bul اینجا مسئله به \key{چهار زیرمسئله} تجزیه می‌شود؛ و تصریح کنید که همین چهارتا با همین ترتیب، ستون‌بندی فصل‌های ۲ تا ۴ است.};

\node[hdr=c1, minimum height=7.4cm] (n3) at (-8.55,0)
 {{\small\bfseries\color{c1} ۱-۳ چالش‌ها}\\[4pt]
  چهار چالش را بگویید و هر کدام را به یک تصمیم روشی در فصل ۳ گره بزنید:\\[3pt]
  \bul \key{فراداده‌ی ناسازگار:} برچسب فاز در سرآیند \lr{DICOM} ناقص است $\leftarrow$ تشخیص فاز از پیش‌نیاز جانبی به مؤلفه‌ی اصلی تبدیل می‌شود.\\[3pt]
  \bul \key{تغییر ناهمگن کنتراست:} عروق و پارانشیم روشن‌تر، استخوان بدون تغییر $\leftarrow$ فرض اطلاعات متقابل می‌شکند.\\[3pt]
  \bul \key{کمیابی جفت طبیعی} $\leftarrow$ هم ارزیابی ثبت و هم آموزش سنتز دشوار می‌شود.\\[3pt]
  \bul \key{ناهمگنی معیار و داده} در ادبیات سنتز $\leftarrow$ انتخاب مبنای معتبر دشوار است.};

\node[hdr=c1, minimum height=6.4cm] (n4) at (8.55,-8.0)
 {{\small\bfseries\color{c1} ۱-۴ اهداف پژوهش}\\[4pt]
  چهار زیرمسئله به چهار هدف عملیاتی ترجمه می‌شوند، به‌علاوه‌ی یک هدف پنجم که اثر تجمیعی را می‌سنجد:\\[3pt]
  ۱. مستندسازی نظام‌مند سوگیری و کفایت مجموعه‌داده‌های عمومی.\\[2pt]
  ۲. مدل تشخیص فاز مستقل از سرآیند.\\[2pt]
  ۳. پایپ‌لاین ثبت با \key{محدودسازی دامنه‌ی} معیار تراز.\\[2pt]
  ۴. آموزش و ارزیابی سنتز روی جفت‌های هم‌ترازشده.\\[2pt]
  ۵. سنجش اثر بر شکاف بخش‌بندی.\\[3pt]
  \bul تصریح کنید: \key{مدل بخش‌بندی بازطراحی نمی‌شود.}};

\node[hdr=c1, minimum height=6.4cm] (n5) at (0,-8.0)
 {{\small\bfseries\color{c1} ۱-۵ ساختار پایان‌نامه}\\[4pt]
  یک پاراگراف نقشه‌ی راه: فصل ۲ پیشینه‌ی پنج حوزه و جای خالی؛ فصل ۳ سه مؤلفه‌ی روش و پروتکل ارزیابی؛ فصل ۴ نتایج تجربی؛ فصل ۵ جمع‌بندی و کارهای آتی.\\[4pt]
  \bul اینجا چیزی برای دفاع نیست؛ فقط خواننده را روی ریل چهار زیرمسئله نگه می‌دارید.};

\node[hdr=c1, minimum height=6.4cm] (n6) at (-8.55,-8.0)
 {{\small\bfseries\color{c1} ۱-۶ جمع‌بندی}\\[4pt]
  \bul پیام واحد فصل: ریشه‌ی افت دقت را نباید فقط در معماری مدل‌ها جست‌وجو کرد.\\[3pt]
  \bul سه عامل \key{مستقل} همگی \key{پیش از} مرحله‌ی مدل‌سازی عمل می‌کنند: کاهش کنتراست بافتی، ناسازگاری فراداده‌ی فاز، و کمیابی داده‌ی جفت‌شده.\\[3pt]
  \bul همین است که مداخله در سطح داده را توجیه می‌کند و پل به فصل ۲ می‌زند.};

\draw[arw] (n1.west) -- (n2.east);
\draw[arw] (n2.west) -- (n3.east);
\draw[arwl] (n3.south) -- ++(0,-0.3) -- ($(n4.north)+(0,0.3)$) -- (n4.north);
\draw[arw] (n4.west) -- (n5.east);
\draw[arw] (n5.west) -- (n6.east);
\end{tikzpicture}
\end{center}

\newpage
% =====================================================================
% PAGE 3 : CHAPTER 2
% =====================================================================
\pagetitle{c2}{فصل ۲ \raisebox{1pt}{---} کارهای پیشین: هر بخش به کدام «جای خالی» می‌رسد}

\setlength{\colw}{5.6cm}
\begin{center}
\begin{tikzpicture}[font=\scriptsize]

\node[hdr=c2, minimum height=7.4cm] (m1) at (9.45,0)
 {{\small\bfseries\color{c2} ۲-۱ پیش‌زمینه‌ی بالینی}\\[4pt]
  \bul تعریف فازهای کنتراست و زمان‌بندی آن‌ها.\\[3pt]
  \bul چرا رادیولوژیست به چند فاز نیاز دارد: الگوی افزایش کنتراست و شست‌وشو.\\[3pt]
  \bul این بخش زمینه است، نه ادعا؛ کوتاه بگویید و فقط واژگانی را بسازید که تا آخر پایان‌نامه لازم دارید (\lr{NC}، \lr{CE}، شریانی، وریدی).};

\node[hdr=c2, minimum height=7.4cm] (m2) at (3.15,0)
 {{\small\bfseries\color{c2} ۲-۲ بخش‌بندی چندعضوی و شکاف کیفیت}\\[4pt]
  \bul اینجا \key{مسئله‌ی چتر} تعریف می‌شود؛ بقیه‌ی فصل زیرمجموعه‌ی آن است.\\[3pt]
  \bul با وجود پیشرفت \lr{nnU-Net} و جانشینانش، افت دقت در فاز بدون‌کنتراست باقی است.\\[3pt]
  \bul دلیل: افت \lr{CNR} و جابه‌جایی حوزه‌ای وابسته به فاز \raisebox{1pt}{---} نه ضعف معماری.\\[3pt]
  \bul نتیجه‌گیری بخش: این افت را باید در داده حل کرد.};

\node[hdr=c2, minimum height=7.4cm] (m3) at (-3.15,0)
 {{\small\bfseries\color{c2} ۲-۳ مجموعه‌داده‌ها و سوگیری}\\[4pt]
  \bul چشم‌انداز مجموعه‌داده‌های عمومی و ابزارهای موجود ارزیابی کیفیت داده.\\[3pt]
  \bul ادعای بخش: هیچ ارزیابی نظام‌مندی از این مجموعه‌داده‌ها \key{به‌عنوان منبع آموزش هوش مصنوعی} وجود ندارد.\\[3pt]
  \bul و هیچ منبعی گزارش نکرده چه سهمی از آن‌ها هر دو فاز را دارند.\\[4pt]
  {\bfseries\color{cred} $\Leftarrow$ جای خالی ۱}};

\node[hdr=c2, minimum height=7.4cm] (m4) at (-9.45,0)
 {{\small\bfseries\color{c2} ۲-۴ تشخیص فاز کنتراست}\\[4pt]
  \bul تعریف مسئله و معیارهای ارزیابی رایج.\\[3pt]
  \bul روش‌های موجود یا شبکه‌ی سه‌بعدی داده‌بر هستند یا به بخش‌بندی از پیش آموزش‌دیده متکی‌اند؛ و به دستگاه و پروتکل حساس‌اند.\\[3pt]
  \bul برچسب سرآیند \lr{DICOM} قابل اتکا نیست، پس این گام پیش‌نیاز است نه اختیاری.\\[4pt]
  {\bfseries\color{cred} $\Leftarrow$ جای خالی ۲}};

\node[hdr=c2, minimum height=7.6cm] (m5) at (9.45,-8.0)
 {{\small\bfseries\color{c2} ۲-۵ ثبت تصویر بین‌فازی}\\[4pt]
  \bul معیارهای کلاسیک (اطلاعات متقابل) بر فرض \key{رابطه‌ی یکنواخت شدت روشنایی در کل تصویر} بنا شده‌اند.\\[3pt]
  \bul تزریق حاجب این فرض را ناهمگن می‌شکند.\\[3pt]
  \bul \key{مهم‌ترین جمله‌ی این بخش:} روش‌های آگاه از آناتومی موجود، آناتومی را برای \key{وزن‌دهی} به کار می‌برند، نه برای \key{محدودسازی دامنه‌ی} معیار.\\[4pt]
  {\bfseries\color{cred} $\Leftarrow$ جای خالی ۳ (سهم اصلی پایان‌نامه از همین‌جا بیرون می‌آید)}};

\node[hdr=c2, minimum height=7.6cm] (m6) at (3.15,-8.0)
 {{\small\bfseries\color{c2} ۲-۶ سنتز \lr{NC\,$\rightarrow$\,CE}}\\[4pt]
  \bul از \lr{GAN}های جفت‌شده تا مدل‌های انتشار اخیر: بازسازی اطلاعات کنتراست ممکن است.\\[3pt]
  \bul اما کیفیت آن به دقت هم‌ترازی هندسی داده‌ی آموزشی وابسته است.\\[3pt]
  \bul تصحیح یک تصور رایج: کد باز در این حوزه به‌کلی غایب نیست.\\[3pt]
  \bul اما هیچ مطالعه‌ای \key{کیفیت هم‌ترازی داده‌ی آموزشی} را گزارش نکرده تا سهم داده از سهم مدل جدا شود.\\[3pt]
  {\bfseries\color{cred} $\Leftarrow$ جای خالی ۴}};

\node[hdr=c2, minimum height=7.6cm] (m7) at (-6.3,-8.0)
 {{\small\bfseries\color{c2} ۲-۷ جمع‌بندی و جای خالی پژوهش}\\[4pt]
  این بخش مهم‌ترین بخش فصل است: چهار جای خالی را کنار هم می‌گذارید و مستقیماً به ساختار فصل ۳ وصل می‌کنید.\\[4pt]
  ۱. تشخیص قابل‌اعتماد فاز، مستقل از سرآیند.\\[2pt]
  ۲. پایپ‌لاینی که \key{دامنه‌ی} محاسبه‌ی معیار تراز را محدود کند، نه وزن آن را.\\[2pt]
  ۳. کمیابی جفت‌های هم‌تراز برای آموزش سنتز.\\[2pt]
  ۴. نبود شواهد کمّی درباره‌ی اثر این مداخلات بر شکاف بخش‌بندی.\\[4pt]
  \bul جمله‌ی پایانی: «فصل بعد دقیقاً به این چهار مورد، به همین ترتیب، می‌پردازد.»};

\draw[arw] (m1.west) -- (m2.east);
\draw[arw] (m2.west) -- (m3.east);
\draw[arw] (m3.west) -- (m4.east);
\draw[arwl] (m4.south) -- ++(0,-0.3) -- ($(m5.north)+(0,0.3)$) -- (m5.north);
\draw[arw] (m5.west) -- (m6.east);
\draw[arw] (m6.west) -- (m7.east);
\end{tikzpicture}
\end{center}

\newpage
% =====================================================================
% PAGE 4 : CHAPTER 3 (part A)
% =====================================================================
\pagetitle{c3}{فصل ۳ \raisebox{1pt}{---} روش پیشنهادی (الف): از نمای کلی تا ثبت تصویر}

\setlength{\colw}{7.7cm}
\begin{center}
\begin{tikzpicture}[font=\footnotesize]

\node[hdr=c3, minimum height=7.3cm] (p1) at (8.55,0)
 {{\small\bfseries\color{c3} ۳-۱ نمای کلی \raisebox{1pt}{---} مهم‌ترین بخش برای دفاع}\\[4pt]
  \bul سه مؤلفه یک \key{زنجیره‌ی وابسته} تشکیل می‌دهند و نه سه کار موازی. اگر داور یک چیز از این فصل ببرد، همین باشد.\\[3pt]
  \bul منطق زنجیره: سنتز به جفت هم‌تراز نیاز دارد $\leftarrow$ جفت باید \key{ساخته} شود $\leftarrow$ ساختن آن دو پیش‌نیاز دارد: شناسایی فاز و هم‌ترازسازی هندسی.\\[3pt]
  \bul \key{نکته‌ی ظریفی که باید تأکید شود:} مدل تشخیص فاز \key{دو نقش} دارد؛ هم دروازه‌ی داده است و هم ابزار اندازه‌گیری در ارزیابی سنتز (بخش ۳-۷).};

\node[hdr=c3, minimum height=7.3cm] (p2) at (0,0)
 {{\small\bfseries\color{c3} ۳-۲ پروتکل مرور نظام‌مند}\\[4pt]
  \bul جست‌وجو در \lr{Google Scholar}، \lr{PubMed}، \lr{Scopus} و \lr{TCIA} به‌علاوه‌ی مخازن غیرمتعارف.\\[3pt]
  \bul پنج معیار ورود (حجمی سه‌بعدی، دارای برچسب، برچسب مرتبط با وظیفه‌ی بالینی، ساختار شکمی، دست‌کم یک مقاله‌ی داوری‌شده).\\[3pt]
  \bul ردیابی منشأ برای کمّی‌کردن \key{بازاستفاده} \raisebox{1pt}{---} گامی که کل یافته‌ی فصل ۴ بر آن استوار است.\\[3pt]
  \bul چارچوب ریسک سوگیری با هفت مقوله و فرایند دو ارزیاب.\\[3pt]
  \bul محدودیت را خودتان بگویید: این مرور \key{چشم‌انداز منتشرشده} را نمایندگی می‌کند، نه کل منابع موجود.};

\node[hdr=c3, minimum height=7.3cm] (p3) at (-8.55,0)
 {{\small\bfseries\color{c3} ۳-۳ داده و پیش‌پردازش}\\[4pt]
  \bul مجموعه‌داده و تقسیم در \key{سطح بیمار} با فایل منجمد \raisebox{1pt}{---} مستقیماً از یافته‌ی نشت داده در فصل ۴ می‌آید.\\[3pt]
  \bul استانداردسازی: بازآرایی به \lr{LPS}، وکسل همسان‌گرد \lr{1.5} میلی‌متر، پنجره‌ی \lr{HU} \key{واحد} در بازه‌ی \lr{$[-200,\,400]$}.\\[3pt]
  \bul چرا پنجره‌ی واحد یک الزام روش‌شناختی است: دو مدل که با دو پنجره سنجیده شوند قابل مقایسه نیستند.\\[3pt]
  \bul \key{تفکیک حیاتی} که تا آخر فصل به آن ارجاع می‌دهید: \key{نقاب ثبت} (دامنه‌ی محاسبه‌ی معیار) در برابر \key{نقاب ارزیابی} (دامنه‌ی سنجش).};

\node[hdr=c3, minimum height=7.6cm] (p4) at (8.55,-7.9)
 {{\small\bfseries\color{c3} ۳-۴ تشخیص فاز کنتراست}\\[4pt]
  \bul \key{فرض طراحی} را صریح بگویید: فاز کنتراست چیزی جز الگوی شدت روشنایی در چند اندام مشخص نیست \raisebox{1pt}{---} رادیولوژیست هم همین کار را می‌کند.\\[3pt]
  \bul پس بازنمایی درست، حجم خام نیست: برای هر اندام یک ویژگی، \key{میانه‌ی} \lr{HU} (میانه و نه میانگین، چون ضایعه و کلسیفیکاسیون داده‌ی پرت‌اند).\\[3pt]
  \bul طبقه‌بند \lr{XGBoost} با اعتبارسنجی متقاطع لایه‌ای؛ دو نسخه روی \lr{WAW-TACE} و \lr{VinDr}.\\[3pt]
  \bul سه مزیت را بشمارید: نیاز به داده‌ی اندک، انتقال‌پذیری میان دستگاه‌ها (ویژگی واحد فیزیکی دارد)، و تفسیرپذیری.};

\node[hdr=c3, minimum height=7.6cm] (p5) at (0,-7.9)
 {{\small\bfseries\color{c3} ۳-۵ الف \raisebox{1pt}{---} ایده‌ی مرکزی ثبت تصویر}\\[4pt]
  \bul اینجا سهم اصلی پایان‌نامه است؛ آهسته و با تأکید بگویید.\\[4pt]
  {\bfseries «معیار تراز نباید در نواحی ناپایدار وزن کمتری بگیرد؛ نباید در آن نواحی \emph{محاسبه} شود.»}\\[4pt]
  \bul تفاوت صوری: $\sum_{x} w(x)s(x)$ در برابر $\sum_{x \in \mathcal{M}} s(x)$.\\[3pt]
  \bul چرا این تفاوت در عمل مهم است: در فرمول وزنی، \key{هیستوگرام مشترک} همچنان از وکسل‌های ناپایدار انباشته می‌شود و چندمُدی می‌ماند. محدودسازی دامنه، هیستوگرام را هم محدود می‌کند.\\[3pt]
  \bul نقاب چهارلایه: استخوان / بافت نرم پایدار / سایر بافت نرم / \key{حذف‌شده} (روده به‌دلیل حرکت دودی و گاز).};

\node[hdr=c3, minimum height=7.6cm] (p6) at (-8.55,-7.9)
 {{\small\bfseries\color{c3} ۳-۵ ب \raisebox{1pt}{---} پایپ‌لاین چندگذره}\\[4pt]
  \bul \key{هم‌ترازی مرکز جرم} با دروازه‌ی اطمینان و دو ناحیه‌ی مرده؛ منطق محافظه‌کاری: یک هم‌ترازی درشتِ نادرست بدتر از نبود آن است.\\[3pt]
  \bul \key{گذر صفر:} حل بسته با \lr{SVD}/کابش روی مراکز جرم اندام‌ها؛ نگهبان دترمینان تا چرخش سره باشد نه بازتاب.\\[3pt]
  \bul \key{گذر یک:} زیان آناتومیک (دایس فرسوده + مرکز جرم) با نلدر-مید \raisebox{1pt}{---} بهینه‌ساز بدون مشتق یک ضرورت است، نه سلیقه.\\[3pt]
  \bul \key{گذر دو:} معیار روی بزرگی گرادیان (کنتراست سطح را جابه‌جا می‌کند، لبه را نه).\\[3pt]
  \bul \key{دروازه‌های پذیرش:} بدترین حالت یک گذر ناموفق «بی‌اثر» می‌شود و نه «مخرب».\\[3pt]
  \bul پالایش نهایی \lr{B}-اسپلاین؛ تاخوردگی گزارش می‌شود اما دروازه نمی‌شود.};

\draw[arw] (p1.west) -- (p2.east);
\draw[arw] (p2.west) -- (p3.east);
\draw[arwl] (p3.south) -- ++(0,-0.25) -- ($(p4.north)+(0,0.25)$) -- (p4.north);
\draw[arw] (p4.west) -- (p5.east);
\draw[arw] (p5.west) -- (p6.east);
\end{tikzpicture}
\end{center}

\newpage
% =====================================================================
% PAGE 5 : CHAPTER 3 (part B)
% =====================================================================
\pagetitle{c3}{فصل ۳ \raisebox{1pt}{---} روش پیشنهادی (ب): سنتز و پروتکل ارزیابی}

\setlength{\colw}{7.7cm}
\begin{center}
\begin{tikzpicture}[font=\footnotesize]

\node[hdr=c3, minimum height=7.5cm] (q1) at (8.55,0)
 {{\small\bfseries\color{c3} ۳-۶ الف \raisebox{1pt}{---} مولد و شرطی‌سازی}\\[4pt]
  \bul صورت‌بندی: ترجمه‌ی تصویر-به-تصویر \key{نظارت‌شده} روی جفت‌های هم‌ترازشده‌ی بخش ۳-۵.\\[3pt]
  \bul \lr{U-Net} دوبعدی روی وصله‌های $128\times128$؛ نرمال‌سازی گروهی (دلیلش در فصل ۴ کمّی می‌شود: مسئله‌ی درز).\\[3pt]
  \bul متمایزکننده‌ی \lr{PatchGAN} به‌صورت \key{شرطی}: نسخه‌ی غیرشرطی فقط واقع‌نمایی را پاداش می‌دهد، نه تطابق با \emph{این} ورودی.\\[3pt]
  \bul زیان: $L_1$ سراسری + $L_1$ وزن‌دار سطح اندام + جمله‌ی متخاصم. روده وزن \key{صفر} در آموزش می‌گیرد اما \key{ارزیابی} می‌شود.\\[3pt]
  \bul مدل \LRE{۲٫۵} بعدی برای زمینه‌ی محور $z$ بدون هزینه‌ی شبکه‌ی سه‌بعدی.\\[3pt]
  \bul شرطی‌سازی سطح \key{یک روش قابل استقرار نیست}؛ اندازه‌گیری یک \key{سقف} است.};

\node[hdr=c3, minimum height=7.5cm] (q2) at (0,0)
 {{\small\bfseries\color{c3} ۳-۶ ب \raisebox{1pt}{---} عدم‌قطعیت و مبناها}\\[4pt]
  \bul استدلال محوری: رگرسیون با زیان $L_1$ در بهترین حالت \key{میانه‌ی شرطی} را می‌دهد. اگر $y=f(x,Z)$ و $Z$ (دوز، زمان‌بندی بولوس، برون‌ده قلبی) مشاهده‌نشده باشد، خطای باقی‌مانده \key{تقلیل‌ناپذیر} است.\\[3pt]
  \bul پس آنچه باید سنجید کاهش خطا نیست، \key{کالیبراسیون} است.\\[3pt]
  \bul دو خانواده: رگرسیون ناهم‌واریانس (مبنای ارزان) و مدل انتشار (که باید هزینه‌اش را توجیه کند).\\[3pt]
  \bul \key{پیش‌بینی ثبت‌شده پیش از آزمایش:} یک نمونه‌ی صادق از توزیع شرطی، با معیار نقطه‌ای حدود $\sqrt{2}$ بدتر امتیاز می‌گیرد ($1.128s$ در برابر $0.798s$). این نقص نیست، بهای صداقت است.\\[3pt]
  \bul مبناهای ادبیات روی \key{همان} تقسیم داده و همان ابزار امتیازدهی بازآموزی می‌شوند \raisebox{1pt}{---} مقایسه‌ی کنترل‌شده، نه بازتولید.};

\node[hdr=c3, minimum height=7.5cm] (q3) at (-8.55,0)
 {{\small\bfseries\color{c3} ۳-۷ الف \raisebox{1pt}{---} چرا معیارهای متعارف کافی نیستند}\\[4pt]
  این خودش یک یافته‌ی روش‌شناختی است، نه مقدمه‌چینی. دو دلیل بگویید:\\[4pt]
  \bul \key{دلیل ساختاری:} پنجره‌ی \lr{HU} هوا، ریه و چربی را به صفر و استخوان را به یک اشباع می‌کند. هر معیاری که روی کل تصویر میانگین بگیرد، عمدتاً روی جرمی میانگین می‌گیرد که همه‌ی مدل‌ها در آن بی‌خطا هستند \raisebox{1pt}{---} حتی یک کپی همانی از تصویر ورودی امتیاز بالایی می‌گیرد.\\[3pt]
  \bul \key{دلیل عمیق‌تر:} \lr{PSNR} و \lr{SSIM} با میانگین شرطی کمینه می‌شوند، پس خروجی \key{تار} را بالاتر از بافت \key{صحیح} امتیاز می‌دهند. یعنی مدل‌ها را در جهت اشتباه مرتب می‌کنند.\\[3pt]
  \bul نتیجه: این معیارها گزارش می‌شوند اما \key{ثانویه}اند.};

\node[hdr=c3, minimum height=7.3cm] (q4) at (8.55,-8.1)
 {{\small\bfseries\color{c3} ۳-۷ ب \raisebox{1pt}{---} معیارهای جایگزین}\\[4pt]
  \bul \key{معیار سرصفحه:} خطای میانه‌ی \lr{HU} در سطح اندام \raisebox{1pt}{---} با همان استخراج ویژگی مدل تشخیص فاز. اینجا زنجیره بسته می‌شود.\\[3pt]
  \bul \key{بازیابی سطح:} شیب رگرسیون $\beta$ (هدف ۱) و نسبت واریانس (هدف ۱). این دو، «مولد» را از «میانگین‌گیر» جدا می‌کنند \raisebox{1pt}{---} مدلی می‌تواند خطای \lr{HU} خوبی داشته باشد و هرگز سطح هیچ بیماری را پیش‌بینی نکرده باشد.\\[3pt]
  \bul \key{بافت و پیوستگی:} چهار نسبت با هدف ۱، نه بیشینه (کمتر از ۱ تاری، بیشتر از ۱ بافت توهمی).\\[3pt]
  \bul \key{کالیبراسیون:} \lr{CRPS} و پوشش بازه‌ی ۹۰ درصد.\\[3pt]
  \bul \lr{LPIPS} و \lr{FID} فقط برای مقایسه‌پذیری با ادبیات، مبنای هیچ ادعایی نیستند.};

\node[hdr=c3, minimum height=7.3cm] (q5) at (0,-8.1)
 {{\small\bfseries\color{c3} ۳-۷ ج \raisebox{1pt}{---} ارزیابی پایین‌دستی و انضباط آماری}\\[4pt]
  \bul تصریح کنید: هیچ‌کدام از معیارهای بالا به پرسش \key{عنوان} پایان‌نامه پاسخ نمی‌دهد. این ارزیابی می‌دهد.\\[3pt]
  \bul مرجع: بخش‌بندی روی تصویر کنتراست‌دار \key{واقعی}. سه بازو: (۱) \lr{NC} واقعی = خودِ شکاف؛ (۲) \lr{CE} تولیدشده؛ (۳) تولیدشده با شرطی‌سازی سطح = سقف.\\[3pt]
  \bul \key{شاهد منفی اسکلتی:} اگر بازوی \lr{NC} روی استخوان هم افت بدهد، آنچه سنجیده‌ایم کیفیت کلی تصویر است نه اثر کنتراست.\\[3pt]
  \bul محدودیت را خودتان بگویید: مرجع، حاشیه‌نویسی انسانی نیست.\\[3pt]
  \bul انضباط آماری: سه بذر و دروازه‌ی $2\sigma$، آزمون‌های زوجی سطح مورد، تفاضل پیکربندی.};

\node[hdr=c3, minimum height=7.3cm] (q6) at (-8.55,-8.1)
 {{\small\bfseries\color{c3} ۳-۸ جمع‌بندی فصل}\\[4pt]
  چهار جمله برای بستن فصل:\\[4pt]
  ۱. مدل تشخیص فاز، مسئله را از حجم خام به بردار کوچکی از میانه‌های \lr{HU} تقلیل می‌دهد و بعداً ابزار سنجش می‌شود.\\[3pt]
  ۲. پایپ‌لاین ثبت، به‌جای وزن‌دهی، \key{دامنه} را محدود می‌کند و با مقداردهی اولیه‌ی بسته از کمینه‌های محلی می‌گریزد.\\[3pt]
  ۳. مدل‌های سنتز همگی روی جفت‌های همین ثبت و یک تقسیم منجمد آموزش می‌بینند.\\[3pt]
  ۴. پروتکل ارزیابی خودش بخشی از سهم است: نشان داده شد چرا معیارهای متعارف قادر به رتبه‌بندی نیستند.\\[4pt]
  \bul جمله‌ی گذار: «فصل بعد نتایج تجربی هر مؤلفه را گزارش می‌کند.»};

\draw[arw] (q1.west) -- (q2.east);
\draw[arw] (q2.west) -- (q3.east);
\draw[arwl] (q3.south) -- ++(0,-0.3) -- ($(q4.north)+(0,0.3)$) -- (q4.north);
\draw[arw] (q4.west) -- (q5.east);
\draw[arw] (q5.west) -- (q6.east);
\end{tikzpicture}
\end{center}

\newpage
% =====================================================================
% PAGE 6 : CHAPTER 4
% =====================================================================
\pagetitle{c4}{فصل ۴ \raisebox{1pt}{---} آزمایش‌ها: چه عددی را با چه ادعایی بگویم}

\setlength{\colw}{7.7cm}
\begin{center}
\begin{tikzpicture}[font=\footnotesize]

\node[hdr=c4, minimum height=7.3cm] (r1) at (8.55,0)
 {{\small\bfseries\color{c4} ۴-۱ الف \raisebox{1pt}{---} ترکیب و افزونگی}\\[4pt]
  \bul ۴۵ مجموعه‌داده، در ظاهر \LRE{۴۷٬۰۴۹} بررسی سی‌تی \raisebox{1pt}{---} اما تنها \LRE{۲۳٬۷۴۶} مورد \key{اصیل}. یعنی \LRE{۴۹٫۵} درصد این اکوسیستم، بازاستفاده است.\\[3pt]
  \bul این بازاستفاده تصادفی نیست: چند مجموعه‌ی پایه (\lr{LiTS}، \lr{KiTS}، \lr{Pancreas-CT}، \lr{BTCV}) گره‌های مرکزی‌اند و به مجموعه‌های ترکیبی سرریز می‌کنند.\\[3pt]
  \bul \key{پیامدی که باید صریح بگویید:} خطر \key{نشت داده} \raisebox{1pt}{---} آموزش روی مجموعه‌ی ترکیبی و ارزیابی روی مجموعه‌ی پایه یعنی احتمالاً همان بیماران در هر دو سو.\\[3pt]
  \bul و همین است که تقسیم در \key{سطح بیمار} در فصل ۳ را توجیه می‌کند.};

\node[hdr=c4, minimum height=7.3cm] (r2) at (0,0)
 {{\small\bfseries\color{c4} ۴-۱ ب \raisebox{1pt}{---} جغرافیا و ریسک سوگیری}\\[4pt]
  \bul با وجود دست‌کم ۲۴ کشور، \LRE{۷۵٫۳} درصد از آمریکا، کانادا و اروپا؛ آمریکا ۲۱ درصد، چین و فرانسه هر یک ۹ درصد.\\[3pt]
  \bul \key{مهم‌تر آنچه غایب است:} آفریقای زیرصحرا، جنوب آسیا و بخش عمده‌ی خاورمیانه هیچ نمایندگی ندارند. مثال بالینی \lr{HCC} در کبد غیرسیروتیک را بگویید تا انتزاعی نماند.\\[3pt]
  \bul روی ۲۲ مجموعه‌داده‌ی بزرگ: ریسک بالای نژادی ۱۶ از ۲۲ و انتخاب ۱۵ از ۲۲، اما برچسب‌گذاری تنها ۵ از ۲۲.\\[3pt]
  \bul \key{تفسیر:} ضعف اصلی این اکوسیستم در \key{کیفیت حاشیه‌نویسی} نیست \raisebox{1pt}{---} که تمرکز رایج پژوهش‌هاست \raisebox{1pt}{---} بلکه در \key{ترکیب} آن است: چه کسانی و چه بیماری‌هایی در داده حاضرند.};

\node[hdr=c4, minimum height=7.3cm] (r3) at (-8.55,0)
 {{\small\bfseries\color{c4} ۴-۱ ج \raisebox{1pt}{---} پیامد برای این پایان‌نامه}\\[4pt]
  سه قید عملی که این یافته‌ها بر کار می‌گذارند:\\[4pt]
  ۱. ارزیابی روی محک‌های ترکیبی بی‌اعتبار است مگر منشأ هر مورد ردیابی شود $\leftarrow$ آزمایش‌ها روی یک مجموعه‌داده‌ی \key{اصیل و منفرد} انجام می‌شوند.\\[3pt]
  ۲. کمیابی داده‌ی جفت‌شده، ثبت و سنتز را از دو گزینه‌ی رقیب به دو \key{ضرورت مکمل} تبدیل می‌کند.\\[3pt]
  ۳. سوگیری جغرافیایی، \key{دامنه‌ی اعتبار} هر نتیجه‌ای را محدود می‌کند؛ این را در تفسیر نتایج صریح نگه دارید.\\[4pt]
  {\bfseries\color{cred} نکته‌ی باز: رقم «چند مجموعه‌داده هر دو فاز را دارند» هنوز از شیت استخراج محاسبه نشده؛ ادعای کمیابی بدون آن پشتوانه‌ی عددی ندارد.}};

\node[hdr=c4, minimum height=7.6cm] (r4) at (8.55,-7.9)
 {{\small\bfseries\color{c4} ۴-۲ الف \raisebox{1pt}{---} قیف داده و اندازه‌ی مسئله}\\[4pt]
  \bul از ۲۶۵ بررسی، تنها ۱۶۹ (\LRE{۶۳٫۸} درصد) سری بدون‌کنتراست دارند. \key{همین عدد}، و نه کیفیت هم‌ترازی، سقف تعداد جفت‌های آموزشی سنتز را تعیین می‌کند.\\[3pt]
  \bul دروازه‌ی سلامت بخش‌بندی ۷۷۱ حجم از ۷۷۲ را پذیرفت؛ تنها حجم مردود هر سه شرط را هم‌زمان نقض می‌کرد \raisebox{1pt}{---} دقیقاً الگویی که دروازه برای آن طراحی شده بود.\\[3pt]
  \bul \key{یافته‌ی مهم و خلاف انتظار:} میانه‌ی ناهم‌ترازی اولیه تنها \LRE{۲٫۷۹} میلی‌متر است و فازها در حدود ۸۷ درصد موارد \key{از پیش هم‌موقعیت}اند.\\[3pt]
  \bul این خلاف فرض رایج ادبیات است و تفسیر کل ابلیشن را عوض می‌کند \raisebox{1pt}{---} آن را به‌عنوان یافته بگویید، نه به‌عنوان عذر.\\[3pt]
  \bul روش گزارش: میانه و چارک (توزیع خطا دُم‌دار است) و آزمون علامت \key{زوجی}.};

\node[hdr=c4, minimum height=7.6cm] (r5) at (0,-7.9)
 {{\small\bfseries\color{c4} ۴-۲ ب \raisebox{1pt}{---} ابلیشن‌ها: آزمون ادعای مرکزی}\\[4pt]
  \bul \key{مهم‌ترین نتیجه‌ی تجربی پایان‌نامه:} محدودسازی دامنه‌ی معیار به نقاب اندام، دقیقاً در همان گروه‌های اندامی که فرضیه پیش‌بینی می‌کرد بهبود معنادار می‌دهد \raisebox{1pt}{---} بافت نرم و عروق، همان‌جا که فرض یکنواختی می‌شکند.\\[3pt]
  \bul و در گروه \key{اسکلتی} \raisebox{1pt}{---} که از پیش به‌عنوان \key{شاهد منفی} معرفی شده بود \raisebox{1pt}{---} سودی نمی‌دهد. این دقیقاً همان چیزی است که فرضیه پیش‌بینی کرده بود، و همین ساختار استدلال را قابل ابطال می‌کند.\\[3pt]
  \bul گذر صفر و گذر یک اثرهای \key{مکمل} و بر گروه‌های متفاوت دارند.\\[3pt]
  \bul دروازه‌های پذیرش، یک گذر ناموفق را از عامل \key{خرابی} به عامل \key{بی‌اثر} تبدیل می‌کنند.};

\node[hdr=c4, minimum height=7.6cm] (r6) at (-8.55,-7.9)
 {{\small\bfseries\color{c4} ۴-۲ ج \raisebox{1pt}{---} مقایسه و صداقت گزارش}\\[4pt]
  \bul \key{خواندن صادقانه‌ی جدول:} \lr{DEEDS} در هر سه معیار دقت بهتر است ($0.886$ در برابر $0.816$) و تفاوت معنادار است. این را اول بگویید، نه آخر.\\[3pt]
  \bul اما \lr{DEEDS} \key{غیرصلب} است و در میانه \LRE{۲٫۹} درصد تاخوردگی دارد؛ روش پیشنهادی صلب است و بنا به ساختار نمی‌تواند تاخوردگی تولید کند. و برای کاربرد پایین‌دستی (ساخت جفت آموزشی)، تاخوردگی بافت را جعل می‌کند.\\[3pt]
  \bul مقایسه‌ی هم‌تراز، صلب با صلب: $0.816$ در برابر $0.761$ برای نسخه‌ی بدون نقاب.\\[3pt]
  \bul نکته‌ی جانبی مفید: \lr{uniGradICON} با بهترین خطای کلی، روی گروه عروقی $5.8$ میلی‌متر خطا دارد در برابر $3.4$ روش پیشنهادی.\\[3pt]
  \bul \key{دو ادعایی که رد شد و باید خودتان بگویید:} هم‌ترازی مرکز جرم دقت نهایی را بهبود نمی‌دهد، و گذر دو در پیکربندی فعلی سودی ندارد.\\[3pt]
  \bul محدودیت: ارقام بر ۲۶ بررسی متناظر استوارند، نه هر ۱۵۶ بررسی ثبت‌شده.};

\draw[arw] (r1.west) -- (r2.east);
\draw[arw] (r2.west) -- (r3.east);
\draw[arwl] (r3.south) -- ++(0,-0.25) -- ($(r4.north)+(0,0.25)$) -- (r4.north);
\draw[arw] (r4.west) -- (r5.east);
\draw[arw] (r5.west) -- (r6.east);
\end{tikzpicture}
\end{center}

\newpage
% =====================================================================
% PAGE 7 : GAPS
% =====================================================================
\pagetitle{cred}{آنچه هنوز گفته نشده \raisebox{1pt}{---} فهرست کارهای باقی‌مانده}

\setlength{\colw}{11.9cm}
\begin{center}
\begin{tikzpicture}[font=\footnotesize]

\node[hdr=cred, minimum height=7.6cm] (s1) at (6.3,0)
 {{\small\bfseries\color{cred} بخش‌های نانوشته}\\[4pt]
  \bul \key{نتایج تجربی تشخیص فاز} (زیرمسئله ۲): روش در بخش ۳-۴ نوشته شده اما هیچ عددی در فصل ۴ نیامده است.\\[3pt]
  \bul \key{نتایج تجربی سنتز} (زیرمسئله ۴): روش در بخش ۳-۶ کامل است؛ نتایج، مقایسه با مبناها و معیارهای کالیبراسیون نوشته نشده‌اند.\\[3pt]
  \bul \key{ارزیابی پایین‌دستی بخش‌بندی}: طرح سه‌بازویی در بخش ۳-۷ ج تعریف شده، اما این تنها سنجه‌ای است که به پرسش عنوان پایان‌نامه پاسخ می‌دهد و هنوز اجرا/گزارش نشده است.\\[3pt]
  \bul \key{فصل ۵ (نتیجه‌گیری)}: به‌کلی خالی است.\\[3pt]
  \bul \key{چکیده‌ی فارسی و انگلیسی}: هنوز متن قالب پیش‌فرض در فایل است.};

\node[hdr=cgy, minimum height=7.6cm] (s2) at (-6.3,0)
 {{\small\bfseries نکات نیازمند تکمیل یا رفع ابهام}\\[4pt]
  \bul \key{رقم فاز کنتراست در مرور}: سهم مجموعه‌داده‌هایی که هر دو فاز را دارند باید مستقیماً از شیت استخراج محاسبه شود. بدون آن، ادعای «کمیابی داده‌ی جفت‌شده» \raisebox{1pt}{---} که پایه‌ی بخش‌های ۳-۵ و ۳-۶ است \raisebox{1pt}{---} پشتوانه‌ی عددی ندارد.\\[3pt]
  \bul \key{هزینه‌ی محاسباتی ثبت تصویر}: زمان اجرا و بیشینه‌ی حافظه‌ی هر گام ثبت نشده و ستون مربوطه در جدول مقایسه خالی است.\\[3pt]
  \bul \key{گسترش مجموعه‌ی ارزیابی ثبت}: از ۲۶ بررسی متناظر به کل ۱۵۶ بررسی ثبت‌شده \raisebox{1pt}{---} در متن به‌عنوان «اولویت نخست» علامت خورده است.\\[3pt]
  \bul \key{شکل ۳-۱}: نمودار جریان کل روش (سه بلوک + ارزیابی پایین‌دستی، با تشخیص فاز در دو نقش) هنوز به‌صورت کامنت در فایل روش است.};

\end{tikzpicture}
\end{center}

\end{document}
